\subsection{Structural extension and independent evaluation}
The extension holds the preparation problem fixed while varying control flow and resource interactions (Table~\ref{tab:processes}). P1 synchronizes two operations sharing a pool. P2 assigns simple and complex cases to different routes and preserves that choice on repetition. P3 combines a local return with repetition of the whole process. P4 has separate arrival streams whose middle operations share capacity. P5 has two bottlenecks that can respond differently to added capacity. P6 combines class-dependent routing, parallel execution and repetition of the selected branch.

\begin{table}[t]\caption{Process families. Counts refer to visible activities, resource pools, distinct return sites and intervention types; no action is available in every family.}\label{tab:processes}\centering\small\setlength{\tabcolsep}{4pt}
\begin{tabular}{lp{.40\textwidth}rrrr}\toprule
 & Structure & Acts. & Pools & Loops & Actions\\\midrule
P0 & Serial sequence and one return & 3 & 3 & 1 & 2\\
P1 & AND block, synchronization, repeat & 4 & 3 & 1 & 2\\
P2 & Class-dependent XOR, branch repeat & 5 & 4 & 2 & 3\\
P3 & Nested local and global repetition & 4 & 4 & 2 & 3\\
P4 & Two processes, shared middle pool & 6 & 5 & 2 & 3\\
P5 & Serial bottlenecks, late-stage return & 5 & 5 & 1 & 3\\
P6 & XOR containing AND, branch repeat & 7 & 6 & 2 & 3\\
\bottomrule\end{tabular}\end{table}

Every new structure is evaluated under fixed, service, routing and combined feedback, with 12 independent training datasets per combination: 288 in total. True service slowdown is a saturating function of queue length; return probability uses a saturating quadratic function of global, local or process-specific load. Each generator, resource assignment, action and cost is fixed before method comparison. Normal-load calibration uses only the generating system and a declared stability grid; it does not consider which method performs best. The archive includes the full specifications and inspected calibration outcomes.

Load thresholds scale with a normal-regime reference. Let $S$ be the rounded-up 95th percentile of WIP in separate fixed-world reference episodes, and $T$ the rounded-up mean case duration. For P1--P6, $S=(3,4,5,5,5,5)$ and $T=(7,7,8,7,10,8)$. Candidate triggers are $2S,4S,8S$, observations are spaced by $0.5T$, and severe overload requires WIP at least $10S$ for $20T$. A separate grid identifies a stable reference arrival rate for each process. Training uses eight episodes at $0.4,0.6,0.8,1.0$ times that rate, twice each; validation uses $0.5,0.7,0.9$. Arrivals end at $40T$, followed by drainage. Every training and validation episode is audited against the severe-event criterion. P0 retains its absolute thresholds and original protocol.

Planning uses six conditions with two repeats and horizon $140T$; testing uses six different conditions with eight repeats and horizon $180T$. The test pulses change both timing and intensity. Separate random streams cover calibration, validation, planning and tests. All candidate actions are tested on the true process after policies are fixed. The known-structure and training-only discovery variants use the same input logs, action catalog and paired test trajectories. Fitness, precision, added or missing paths and executable-model failures are recorded without repairing discovered structures.

Besides static, predictive-selected and retained-load models, the extension evaluates service-only and routing-only ablations, individual return-site ablations, every fixed action/threshold, and a known-generator planner. A feedback-strength rule selects an action from estimated local service or repetition amplification and uses threshold $2S$; its inputs come from training data, without simulating alternatives. Generator-planner regret compares realized costs of policies selected using the same finite planning library. A test-selected minimum is reported only as a hindsight bound.

The primary new contrasts compare retained-load with predictive selection under routing feedback, separately for P1--P6 with known structure. We average trajectories within each trained model and use paired differences over the 12 independent fits. Bonferroni correction covers the six primary contrasts; other intervals are descriptive. Cross-process summaries give equal weight to each process after normalization by its no-action cost. Known and discovered structures remain separate. Diagnostic tests vary selector thresholds, shift P2's complex-case fraction from $0.35$ to $0.65$, and change P3's actual rework reduction to $30\%$ or $70\%$ while keeping nominally selected policies fixed. These tests do not retune the method.
